\chapter{Genre: The Rarely/Barely Pair}\label{ch:genre}

The idea behind this book began with a pair of films. \emph{Rarely Needed Protocols} was to be a science-fiction thriller: a young astronaut named Theseus arrives on a planet in civil war, discovers that its ruling class controls the population through technology, sabotages the control systems, and must live with the chaos his sabotage unleashes. \emph{Barely Needed Protocols} was to be a satirical comedy for children about the same world and the same events, mocking the absurdity of a regime that bans square televisions and permits only round ones. The two were to be screened at the same time, on the same screen, separated by the timing of the viewers' glasses, and joined by a soundtrack written to serve both.

This chapter takes the pair as the founding case of genre variation. It argues that genre is the most natural dimension for a variant family, and the one that most clearly shows why such a family has to be written together from the start.

\section{Complementary, not restrictive}

It is worth being precise about what the pair was and was not. It was not an adult film with a children's cut, in which material is removed so that a younger audience can safely see what remains. \Cref{ch:unfiltered} showed that the glasses cannot enforce such a removal: a child without glasses sees the composite, and the adult material is in it. The pair was two complete tellings of one story, each addressed to its own audience, neither containing anything the other audience needed to be protected from.

The distinction is not only technical. A restrictive family is organized around what one audience must not see, and its variants are defined by absence. A complementary family is organized around what each audience is offered, and its variants are defined by address. The thriller treats the regime's control as a horror and Theseus's sabotage as a moral wager with real costs. The satire treats the same control as ridiculous and the same sabotage as a prank that goes gloriously wrong. Both are true to the events. Neither is a diminished form of the other. A parent and child who watch them side by side have seen the same story told to each of them.

\section{Distances}

In the vocabulary of \Cref{ch:family}, a genre pair has a distinctive profile.

Its fact distance is small by design. Theseus arrives, discovers the control systems, sabotages them, and faces the consequences in both realizations. What is established about the world and its events is the same, even when the reasons given for them differ in emphasis. The realization graph therefore stays close to complete, and a viewer can move between the thriller and the satire at almost any point without losing their place in the story. This is the strongest argument for genre as a first dimension: it gives the richest switching of any kind of variation.

Its image distance is large. Genre lives in performance, framing, lighting, and cutting, and all of these are on screen. Unless the pair shares footage, which is possible to a limited degree, the realizations show different pictures for most of the running time, and the costs of multiplexing in light and flash rate, examined in \Cref{ch:capacity}, are paid throughout.

Its time distance is where the work lies.

\section{Genre keeps time}

Comedy and drama do not only look different. They move differently. A joke depends on timing measured in fractions of a second: the pause before a reaction, the beat before a repetition, the interruption that arrives exactly too early. Drama often depends on the opposite, on holding a shot longer than is comfortable, letting silence do the work. The same event, played for laughs and played for weight, rarely takes the same number of seconds.

A variant family cannot require its realizations to match second by second, and it does not need to. It needs them to match at the points where it matters, and it can let them diverge in between. The natural unit is the scene. Assign each scene $j$ an \emph{envelope}, an interval $[s_j,e_j]$ on the common clock, and require every realization to begin and end the scene within it. Inside the envelope each realization keeps its own time. Write $\delta_k(t)$ for how far realization $k$ is ahead of or behind the scaffold at time $t$. The envelope condition says
\begin{equation}
\delta_k(s_j)=\delta_k(e_j)=0\quad\text{for every scene } j,
\end{equation}
while $\delta_k(t)$ may vary freely within the scene.

The envelope turns genre timing from an obstacle into a constraint that can be written to. A comic realization that needs three extra seconds for a reaction must give them back before the scene ends, perhaps by cutting a line the drama needs and the comedy does not. A dramatic realization that holds on a face must find the time in a transition the comedy rushes through. The envelope boundaries are also the natural candidates for structural rendezvous, because both realizations have completed the same action there, and the realization graph is most reliably complete at the moments the envelopes enforce.

\section{Three ways to make a genre pair}

There are three production strategies, and they trade image distance against cost and freedom.

The first uses the same footage for both realizations and varies only editing, music, and sound. This is the cheapest and the most constrained. Image distance is small, because many frames coincide, so the light cost is low and adaptive gating opens the glasses much of the time. But a comedy cut from dramatic footage, or the reverse, is limited by performances that were given for one purpose. Editing can do a surprising amount, as trailers that recut a drama as a comedy demonstrate, but it cannot make a solemn actor look mischievous.

The second shoots the same blocking twice with different performances. The actors hit the same marks, the camera takes the same positions, the action happens in the same order, and the delivery differs: one take for the thriller, one for the satire. Image distance is moderate. Many shots coincide in composition and differ in expression. The strategy preserves the structural correspondence that keeps envelopes easy to meet, and it gives each realization its own performances.

The third stages the two realizations separately, each shot in its own way, joined only at the envelope boundaries. It gives the most freedom and costs the most, in production and in light, because the realizations share almost no frames. It is appropriate when the genres want genuinely different visual languages, which a thriller and a children's satire plausibly do.

Nothing requires one strategy throughout. A pair might share establishing shots and action, reshoot dialogue scenes with different performances, and stage a few set pieces separately. The choice can be made scene by scene, with the image budget in view.

\section{Genre is written in}

The strategies share a premise that is easy to miss. In each of them, the story must already be able to bear both readings. A thriller cannot be made into a comedy after the fact without a comedy's structure, and a satire for children cannot be turned into a thriller without a thriller's stakes. The events that both realizations share have to be events that are both funny and serious, or at least capable of being told as either. This is a constraint on the story before it is a constraint on the film.

It is also a real artistic opportunity. Stories that support both readings are not rare; much of the best comedy is about serious things, and much of the best drama has absurdity in it. The \emph{Protocols} pair is an instance. A regime that controls its population by controlling the shape of its televisions is sinister and ridiculous at once, and each realization develops one side of a situation that already contains both. The strongest genre pairs will be those whose shared events contain the tension their realizations divide between them.

There is a characteristic difficulty. Comedy and drama sometimes disagree about what must be shown. A comic beat may depend on a visual detail the drama has no reason to include, and a revelation in the drama may need a close-up the comedy would find heavy. Where the difference is only in emphasis, both realizations can show the detail and weight it differently. Where one realization needs to establish something the other does not, the difference becomes fact distance, and the pair must either reconcile it at the next envelope boundary or accept that switching is closed until it does. The grammar of \Cref{ch:divergence} predicts exactly this: genre variation is cheap in fact distance only as long as the genres agree on what happens.

\section{The shared score}

The original conception of the pair gave it one soundtrack, composed so that it would serve both films at once. That choice is examined fully in \Cref{ch:composite-sound}, but it bears on genre directly. Music is one of the strongest carriers of genre, and a score that serves a thriller and a satire simultaneously cannot simply be thrilling or funny. It must be ambiguous in the way the shared events are ambiguous, supporting tension in one picture and absurdity in the other. That is difficult, and it is also a familiar achievement: film scores often carry irony, and a single cue can read as menace or mock-menace depending on what it accompanies. The shared score makes the genre pair a single work in the room even while the pictures diverge.

\section{The composite of two genres}

A final possibility follows from \Cref{ch:unfiltered}. A viewer without glasses, or with a coarser setting that admits both streams, sees the average of the thriller and the satire. Ordinarily that average is a double exposure of two performances. But a pair made with the second strategy, the same blocking with different delivery, produces a composite in which the same figures stand in the same places with two expressions superimposed. That composite is neither thriller nor comedy. It is something closer to a tragicomic ghost image, and it may be watchable in its own right as the room's shared view of an event that is both. Whether it is worth designing for is an artistic question. That it exists at all, as a third realization no one had to write, is a property only this medium has.
